% 原书第 22—26 页；第 2.2 节及图 2.3。
\section{历史札记}
\label{sec:historical-notes}

自主性和不变性的思想，深深植根于\textbf{结构方程模型（structural equation model，SEM）}或结构因果模型的概念之中。我们更倾向于采用后一个术语，因为在不少使用 SEM 这一名称的语境中，人们把结构赋值当作代数方程，而不是赋值。相关文献十分广泛，可参阅 Aldrich［1989］、Hoover［2008］和 Pearl［2009］的综述。

结构方程模型的一个思想先驱，是 Wright［1918, 1920, 1921］开创的\textbf{路径模型（path model）}概念（见图 2.3）。尽管 Wright 是生物学家，如今结构方程模型与计量经济学的联系最为紧密。据 Hoover［2008］记述，Jan Tinbergen 在 20 世纪 30 年代开创了结构计量经济模型的研究，而概率计量经济学的概念基础则由 Trgyve Haavelmo 的工作奠定［Haavelmo, 1944］。早期经济学家试图从概念上阐明一个事实：与相关不同，回归具有自然的方向性。$Y$ 对 $X$ 回归所得的解，通常不是 $X$ 对 $Y$ 回归所得解的逆。\footnote{顺便提及，尽管早期工作大多只使用线性方程，也有人尝试将其推广到非线性结构方程模型［Hoover, 2008］。}
那么，数据究竟如何告诉我们，应该沿哪个方向进行回归？这是\textbf{观测等价性（observational equivalence）}问题，与计量经济学家所谓的\textbf{识别（identification）}问题密切相关。

\begin{figure}[p]
\centering
\begin{tikzpicture}[causal,x=1.05cm,y=.9cm,
 ch02dot/.style={circle,fill=black,inner sep=1.8pt},
 ch02coef/.style={font=\small,fill=white,inner sep=.6pt},
 ch02pigbody/.pic={
   \path[draw,fill=white,line width=.7pt]
   (-.95,-.36) .. controls (-1.15,.1) and (-.75,.55) .. (-.3,.44)
   .. controls (.1,.37) and (.33,.63) .. (.61,.45)
   .. controls (.92,.2) and (1.13,.13) .. (.94,-.04)
   --(.73,-.17)--(.77,-.34)--(.28,-.34)
   .. controls (-.15,-.39) and (-.6,-.27) .. cycle;
   \fill (.76,.21) circle (1pt);
 }]
% 原图中的四只豚鼠，及各个体遗传构成和生殖细胞遗传构成。
\begin{scope}[shift={(3.65,2.25)}]
\pic {ch02pigbody};
\path[draw,pattern=north west lines] (-.3,.44) .. controls (-.77,.56) and (-1.03,.04) .. (-.95,-.25) .. controls (-.7,-.38) and (-.35,-.37) .. (-.23,-.26)--cycle;
\path[draw,pattern=north west lines] (.5,.45) .. controls (.25,.22) and (.42,-.13) .. (.88,.02) .. controls (.88,.18) and (.64,.4) .. cycle;
\end{scope}
\node[font=\small] at (3.6,2.98){父本};
\begin{scope}[shift={(3.65,-6.25)}]
\pic {ch02pigbody};
\path[draw,pattern=north west lines] (-.55,.38) .. controls (-.77,.42) and (-.81,-.12) .. (-.52,-.26) .. controls (-.2,-.4) and (.13,-.3) .. (.05,.31)--cycle;
\path[draw,pattern=north west lines] (.82,.29)--(.72,.35)--(.77,-.12)--(.96,-.03)--cycle;
\end{scope}
\node[font=\small] at (3.6,-7.02){母本};
\begin{scope}[shift={(10.25,0)}]
\pic {ch02pigbody};
\path[draw,pattern=north west lines] (-.55,.38) .. controls (-.77,.42) and (-.81,-.12) .. (-.52,-.26) .. controls (-.2,-.4) and (-.09,-.03) .. (-.19,.34)--cycle;
\path[draw,pattern=north west lines] (.5,.45) .. controls (.25,.22) and (.42,-.13) .. (.88,.02) .. controls (.88,.18) and (.64,.4) .. cycle;
\end{scope}
\node[font=\small] at (10.25,.78){$O$};
\begin{scope}[shift={(10.25,-4.3)}]
\pic {ch02pigbody};
\path[draw,pattern=north west lines] (-.08,.43) .. controls (-.37,.45) and (-.42,-.1) .. (-.12,-.16) .. controls (.25,-.23) and (.34,-.01) .. (.2,.45)--cycle;
\end{scope}
\node[font=\small] at (10.25,-3.52){$O'$};
\node[ch02dot,label=above left:$H''$] (ch02hpp) at (1.45,0) {};
\node[ch02dot,label=above left:$H'''$] (ch02hppp) at (1.45,-4.3) {};
\node[ch02dot,label=above:$G$] (ch02g) at (4.5,0) {};
\node[ch02dot,label=right:$G'$] (ch02gp) at (4.5,-1.08) {};
\node[ch02dot,label=right:$G''$] (ch02gpp) at (4.5,-3.22) {};
\node[ch02dot,label=below:$G'''$] (ch02gppp) at (4.5,-4.3) {};
\draw (4.48,.06) .. controls (4.25,.14) and (4.46,.3) .. (4.25,.3) .. controls (4.10,.28) and (4.24,.44) .. (4.06,.45);
\draw (4.48,-1.02) .. controls (4.25,-.94) and (4.46,-.78) .. (4.25,-.78) .. controls (4.10,-.8) and (4.24,-.64) .. (4.06,-.63);
\node[ch02dot,label=above right:$H$] (ch02h) at (7.5,0) {};
\node[ch02dot,label=above right:$H'$] (ch02hp) at (7.5,-4.3) {};
\node[font=\small] (ch02chance) at (.6,-2.1){随机性};
\node (ch02sd) at (.25,2.7){$D$};
\node (ch02se) at (.25,1.65){$E$};
\node (ch02me) at (.25,-5.83){$E$};
\node (ch02md) at (.25,-6.88){$D$};
\node (ch02od) at (8.8,2.15){$D$};
\node (ch02opd) at (9.18,-6.45){$D$};
\node (ch02oe) at (7.5,-6.45){$E$};
\draw[ed] (ch02sd)--node[ch02coef,above]{$d$}(2.55,2.56);
\draw[ed] (ch02se)--node[ch02coef,above]{$e$}(2.55,2.00);
\draw[ed] (ch02hpp)--node[ch02coef,left]{$h$}(2.65,1.7);
\draw[ed] (ch02me)--node[ch02coef,above]{$e$}(2.55,-6.00);
\draw[ed] (ch02md)--node[ch02coef,above]{$d$}(2.55,-6.55);
\draw[ed] (ch02hppp)--node[ch02coef,left]{$h$}(2.65,-5.70);
\draw[ed] (ch02hpp)--node[ch02coef,above,pos=.7]{$b$}(ch02g);
\draw[ed] (ch02hpp)--node[ch02coef,above,pos=.3]{$b$}(ch02gp);
\draw[ed] (ch02hppp)--node[ch02coef,below,pos=.3]{$b$}(ch02gpp);
\draw[ed] (ch02hppp)--node[ch02coef,above,pos=.7]{$b$}(ch02gppp);
\draw[ed] (ch02chance.east)--node[ch02coef,above,pos=.55]{$c$}(ch02g);
\draw[ed] (ch02chance.east)--node[ch02coef,below,pos=.7]{$c$}(ch02gp);
\draw[ed] (ch02chance.east)--node[ch02coef,above,pos=.55]{$c$}(ch02gpp);
\draw[ed] (ch02chance.east)--node[ch02coef,below,pos=.7]{$c$}(ch02gppp);
\draw[ed] (ch02g)--node[ch02coef,above]{$a$}(ch02h);
\draw[ed] (ch02gpp)--node[ch02coef,above,sloped,pos=.65]{$a$}(ch02h);
\draw[ed] (ch02gp)--node[ch02coef,above,sloped,pos=.65]{$a$}(ch02hp);
\draw[ed] (ch02gppp)--node[ch02coef,above]{$a$}(ch02hp);
\draw[ed] (ch02h)--node[ch02coef,above]{$h$}(9.12,0);
\draw[ed] (ch02hp)--node[ch02coef,above]{$h$}(9.12,-4.3);
\draw[ed] (ch02od)--node[ch02coef,above left]{$d$}(9.65,.65);
\draw[ed] (ch02opd)--node[ch02coef,above left]{$d$}(9.78,-4.83);
\draw[ed] (ch02oe)--node[ch02coef,left,pos=.6]{$e$}(9.5,-.65);
\draw[ed] (ch02oe)--node[ch02coef,below right,pos=.6]{$e$}(9.3,-4.85);
\end{tikzpicture}
\par\smallskip
\begin{minipage}{.96\textwidth}
\small
\textbf{图内说明（原图 5）：}该图展示了同窝个体（$O,O'$）之间，以及它们各自与父母之间的因果关系。$H,H',H'',H'''$ 表示四个个体的遗传构成，$G,G',G'',G'''$ 表示四个生殖细胞的遗传构成。$E$ 表示同窝个体共有的环境因素。$D$ 表示其他因素，主要是个体发育中的不规则变异。小写字母表示各条路径的系数。
\end{minipage}
\caption{早期路径图。其中，母本（dam）和父本（sire）分别指豚鼠的雌性亲本和雄性亲本。路径系数刻画给定路径的重要程度，其定义如下：除所考察原因之外，其他所有原因均保持恒定，而所考察原因自身的变异程度保持不变，此时结果的变异程度与其总变异程度之比。依据 Wright［1920］原图重绘。}
\label{fig:ch02-wright-path}
\end{figure}


一些早期工作已经注意到，使一组方程或关系具有“结构性”的特征［Frisch and Waugh, 1933］，与不变性和自主性之间存在联系。按照 Aldrich［1989］的说法，这恰恰是 Frisch et al.［1948］开创性工作的核心概念。这里，结构关系所追求的，不只是对观测到的数据分布进行建模，而是刻画连接模型中各个变量的底层结构。

当时，Cowles 委员会是一家重要的经济研究机构，对计量经济学这一领域的形成起到了关键作用。该机构的工作把因果性与结构计量经济模型的不变性联系起来［Hoover, 2008］。Pearl［2009］认为，“结构方程在系统发生某些变化时保持不变”这一思想，可追溯至 Marschak 为 1950 年 Cowles 专著撰写的开篇章节［Marschak, 1950］。Cowles 的工作着重强调\textbf{内生变量（endogenous variable）}与\textbf{外生变量（exogenous variable）}的关键区别。内生变量是建模者试图理解的变量；外生变量则由模型之外的因素决定，并被视为给定。Koopmans［1950］考察了判断哪些变量应作为外生变量的两项原则。“学科划分原则”把学科范围之外的变量视为外生变量，例如，天气相对于经济学是外生的；更受推崇的“因果原则”，则把那些影响其余内生变量、却几乎不受这些变量影响的变量称为外生变量。

Haavelmo［1943］将结构方程解释为有关假想受控实验的陈述。他考察了带环的随机方程模型，并讨论了不变性的作用以及政策干预。Pearl［2015］评价了 Haavelmo 在政策干预问题研究和因果推断领域发展中的贡献。Hoover［2008］在回顾经济学和计量经济学中的因果性时，讨论了如下形式的系统：
\[
 \begin{aligned}
 X^i&\coloneqq N_X^i,\\
 Y^i&\coloneqq \theta X^i+N_Y^i,
 \end{aligned}
\]
其中，误差 $N_X^i,N_Y^i$ 独立同分布，$\theta$ 是参数。他将下述不要求任何时间顺序的观点归于 Simon［1953］：可以说 $X^i$ 导致 $Y^i$，因为即使不了解 $Y^i$，我们也能完全了解 $X^i$，反之却不成立。这些方程也使我们能够预测干预的作用。Hoover 进一步指出，可以交换 $X^i$ 和 $Y^i$ 的角色，重写这个系统，同时仍使误差项保持不相关。\footnote{第 4.1.3 节将更详细地重新讨论这一问题。}
因此，仅凭一组数据，我们无法推断正确的因果方向，这就是“观测等价性”。受控实验或自然实验则可能帮助我们作出判断。例如，如果某项实验能够改变给定 $X^i$ 时 $Y^i$ 的条件分布，却不改变 $X^i$ 的边缘分布，那么必然是 $X^i$ 导致 $Y^i$。Hoover 将其称为 Simon 的不变性准则：真正的因果顺序，就是在恰当类型的干预下保持不变的顺序。\footnote{我们认为，如果干预之间存在耦合，例如，若干干预相互配合，使反因果条件分布（给定结果时原因的条件分布）保持不变，这一说法可能不成立。可以将这种情形看作在干预层面违反了原理 2.1。第 2.3.4 节将重新讨论这一点。}

Hurwicz［1962］指出，一个方程系统之所以具有结构性，是因为它在某一类修改下保持不变。这样的系统因而与自然定律具有相似之处。Hurwicz 认识到，可以利用这些修改来确定结构；他还指出，因果性离不开结构，而预测却不一定需要结构。

Aldrich［1989］讨论了自主性在结构方程建模中的作用。他认为，自主关系可能比其他关系更稳定，并将 Haavelmo 的自主变量等同于后来所谓的外生变量。自主变量是由外部力量固定，或在处理时被视为随机独立的参数。\footnote{这类似于我们在结构因果模型中采用的噪声项独立假设。}
据 Aldrich［1989，第 30 页］所述，“使用‘自主’这一限定语，以及‘所考察部门之外的力量’这一说法，意味着……该模型的参数不会随部门参数的变化而变化。”他还将不变性与一种称为\textbf{超外生性（super-exogeneity）}的概念联系起来［Engle et al., 1983］。

尽管结构方程建模的早期倡导者已经对其因果基础形成了某些深刻认识，计算机科学中的相关发展最初却独立展开。Pearl［2009，第 104 页］叙述了他和合作者如何开始把贝叶斯网络与结构方程建模联系起来：
\begin{quote}
“我突然意识到，经济学家与统计学家之间长达一个世纪的紧张关系，源于一个简单的语义混淆：统计学家把结构方程理解为关于 $\EE[Y\mid x]$ 的陈述，而经济学家把它们理解为关于 $\EE[Y\mid \dox{x}]$ 的陈述。这就解释了，为什么统计学家声称结构方程没有意义，而经济学家反驳说统计学缺乏实质内容。”
\end{quote}
Pearl［2009，第 22 页］将独立性原理表述为：“网络中的每一种父子关系都代表一种稳定而自主的物理机制——换言之，可以设想，在改变某一关系时，其他关系保持不变。”

值得注意的是，在图 2.2 所示原理 2.1 的不同推论中，使用因果贝叶斯网络的大多数工作，仅仅利用了噪声项的独立性。\footnote{某些贝叶斯结构学习方法［Heckerman et al., 1999］，为给定原因时每个变量的条件概率赋予相互独立的先验，可以将其看作对独立性原理的一种实现。}
这也是撰写本书的动机之一。噪声项的独立性导出了丰富的条件独立结构［Pearl, 2009；Spirtes et al., 2000；Dawid, 1979；Spohn, 1980］，其根源最终可追溯到 Reichenbach 的原理 1.1。相比之下，独立性的其他方面受到的关注明显较少［Hausman and Woodward, 1999；Lemeire and Dirkx, 2006］，但近来已有一系列工作尝试对它们进行形式化并加以利用。其中一个主要动机是原因—结果问题：在这一问题中只有两个变量，因此条件独立性派不上用场（见第 4.1.2 节和第 6.10 节）。

Janzing and Schölkopf［2010］用\textbf{算法信息论（algorithmic information theory）}来形式化机制的独立性（见第 4.1.9 节）。他们认为，结构因果模型中的函数代表独立的因果机制；即使操纵输入的分布或其他机制，这些机制仍然保持不变。更具体地说，在因果贝叶斯网络中，他们假定，给定各节点的父节点时，所有节点的条件分布在算法意义上相互独立。特别地，对于因果贝叶斯网络 $X\to Y$，$P_X$ 与 $P_{Y\mid X}$ 不包含有关彼此的算法信息；也就是说，知道其中一个，并不能使另一个的描述变得更短。把结构因果模型从随机变量推广到个别对象，并以算法相依取代统计相依，就会得到“互无关联的机制在算法意义上相互独立”这一思想。

Schölkopf et al.［2012，例如第 2.1.1 节］讨论了在两个变量的情形下，对原因分布发生变化的稳健性问题，并将其与机器学习问题联系起来，亦见第 5 章。在结构因果模型框架内，他们针对不同的学习情景（如迁移学习、概念漂移），分析函数或噪声的不变性。他们所使用的“机制与输入的独立性”概念，同时涵盖了变化下的独立性和信息论意义上的独立性；在讨论图 2.2 的各个方框时，我们将这种联系称为第一种与第二种独立性之间的“重叠”：“$P_{E\mid C}$ 不包含有关 $P_C$ 的信息，反之亦然；特别地，如果 $P_{E\mid C}$ 在某个时刻发生变化，我们没有理由认为 $P_C$ 也会同时变化。”

关于它们与迁移学习及相关机器学习问题的进一步联系，可参阅 Bareinboim and Pearl［2016］、Rojas-Carulla et al.［2016］、Zhang et al.［2013］和 Zhang et al.［2015］。Peters et al.［2016］利用跨环境的不变性，学习了多元结构因果模型背后图结构的一部分（见第 7.1.6 节）。
\FloatBarrier
